\chapter{有效作用量与费米液体}
\renewcommand{\currentchapterpath}{chapters/ch07/}

\begin{flushright}
\textit{“执本而末自从。” ——《淮南子》
\footnote{有效作用量抓住不可约骨架；Landau 费米液体则抓住低能准粒子这一“本”。}}
\end{flushright}

有效作用量把连通关联的生成泛函经 Legendre 变换变为只含不可约图的骨架展开；Luttinger--Ward 泛函给出自能与热力学的统一框架。Landau 费米液体理论则是金属低能物理的现象学与微观基础。有效作用量部分对齐施均仁讲义第 5 章；费米液体部分对齐 Giuliani--Vignale 第 8 章与 Gross 第 30--34 章（含微观导出与 Kondo）。


\chapterinput{effective_action}
\chapterinput{luttinger_ward}
\chapterinput{fermi_liquid}
\chapterinput{landau_thermo}
\chapterinput{landau_transport}
\chapterinput{landau_micro}
\chapterinput{kondo}
\chapterinput{generalizations}
